\documentclass[11pt,a4paper]{article}

% ── Packages ────────────────────────────────────────────────────────────
\usepackage[utf8]{inputenc}
\usepackage[T1]{fontenc}
\usepackage{amsmath,amssymb,amsfonts,amsthm}
\usepackage{mathtools}
\usepackage{booktabs}
\usepackage{graphicx}
\usepackage{geometry}
\usepackage{hyperref}
\usepackage{cleveref}
\usepackage{listings}
\usepackage{xcolor}
\usepackage{enumitem}
\usepackage{array}
\usepackage{tabularx}
\usepackage{multirow}
\usepackage{caption}
\usepackage{subcaption}
\usepackage{algorithm}
\usepackage{algpseudocode}

\geometry{margin=2.5cm}
\hypersetup{colorlinks=true,linkcolor=blue!70!black,citecolor=green!50!black,urlcolor=blue!70!black}

% ── Theorem environments ──────────────────────────────────────────────────
\newtheorem{theorem}{Theorem}
\newtheorem{lemma}{Lemma}
\newtheorem{proposition}{Proposition}
\newtheorem{corollary}{Corollary}
\theoremstyle{definition}
\newtheorem{definition}{Definition}
\newtheorem{example}{Example}
\theoremstyle{remark}
\newtheorem{remark}{Remark}

% ── Custom commands ──────────────────────────────────────────────────────
\newcommand{\R}{\mathbb{R}}
\newcommand{\C}{\mathbb{C}}
\newcommand{\N}{\mathbb{N}}
\newcommand{\Z}{\mathbb{Z}}
\newcommand{\norm}[1]{\left\lVert#1\right\rVert}
\newcommand{\abs}[1]{\left\lvert#1\right\rvert}
\newcommand{\conj}[1]{\overline{#1}}
\newcommand{\Real}{\operatorname{Re}}
\newcommand{\Imag}{\operatorname{Im}}
\newcommand{\sign}{\operatorname{sgn}}
\newcommand{\FFT}{\operatorname{FFT}}
\newcommand{\IFFT}{\operatorname{IFFT}}
\newcommand{\DFT}{\operatorname{DFT}}
\newcommand{\softthr}{\operatorname{soft\_thr}}
\newcommand{\modReLU}{\operatorname{modReLU}}
\newcommand{\bigO}{\mathcal{O}}

% ── Title ────────────────────────────────────────────────────────────────
\title{\textbf{Spectral Silicon: A Custom ASIC for LLM Inference via\\
Neural-Operator-Based Spectral Mixing}\\[0.5em]
\large Technical Design Document}
\author{Walker Kirkpatrick}
\date{August 2026}

\begin{document}
\maketitle

\begin{abstract}
This document presents the complete architecture, mathematical foundations, and
performance analysis of \emph{Spectral Silicon}, a custom ASIC that replaces the
$\bigO(n^2)$ self-attention mechanism in large language models with an $\bigO(n \log n)$
spectral mixing pipeline based on Fourier Neural Operators (FNO), Adaptive FNO (AFNO),
and FFTNet. The chip implements the pipeline \textsc{FFT} $\to$ spectral weight multiply
$\to$ \textsc{IFFT} $\to$ modReLU directly in silicon on the open-source SkyWater
SKY130 130\,nm process, targeting fabrication through Tiny Tapeout (\$50--\$500).
The design comprises 84 Verilog RTL modules, 12 Python simulation modules, and 343
passing tests. Three architectural generations (V1--V3) are analyzed: V1 achieves
189{,}000 tokens/s at 100{,}300 gate equivalents (GE); V2 reduces area by 91\%
(9{,}102 GE) through a shared FFT/IFFT engine and serialized channels; V3 adds 20
performance improvements including Booth multipliers, block floating-point, and
dual-channel processing, reaching 21{,}900 tokens/s at 80\,MHz (120\,MHz in V6/V7).
All 11 security measures---including constant-time MAC, power flattening, and
bitstream encryption---are preserved across all versions.
\end{abstract}

\tableofcontents
\newpage

%═══════════════════════════════════════════════════════════════════════════
\section{Introduction and Motivation}
%═══════════════════════════════════════════════════════════════════════════

\subsection{The Attention Bottleneck}

Modern large language models (LLMs) are bottlenecked by \emph{self-attention}---the
$\bigO(n^2)$ operation where every token attends to every other token. For a
sequence of length $n=256$ with $d=64$ channels, the attention mechanism computes:
\begin{equation}
  \text{Attention}(Q, K, V) = \text{softmax}\!\left(\frac{QK^\top}{\sqrt{d_k}}\right) V
  \label{eq:attention}
\end{equation}
where $Q, K, V \in \R^{n \times d}$ are the query, key, and value projections. The
matrix $QK^\top \in \R^{n \times n}$ requires $n^2 d$ multiply-accumulate operations
and $n^2$ storage. At 130\,nm, a $256 \times 256$ attention matrix needs
$\sim$256\,KB of SRAM and 65{,}536 MACs---far exceeding the area available on a
Tiny Tapeout tile ($\sim$1\,mm$^2$, $\sim$100K gates).

\subsection{The Spectral Alternative}

Neural operators [?] offer an
alternative: replace the $\bigO(n^2)$ attention matrix with an $\bigO(n \log n)$
spectral mixing operation that performs the \emph{same job} (token mixing) but
through the Fourier domain:
\begin{equation}
  y = x + \IFFT\!\left( R \cdot \FFT(x) \right)
  \label{eq:spectral_mixing}
\end{equation}
where $R$ is a learnable complex weight tensor applied to the first $k \ll n$ Fourier
modes. The key insight is that \textbf{the weights live in the frequency domain},
indexed by mode rather than by position, so the parameter count is $\bigO(k \cdot d)$
instead of $\bigO(n^2)$, and the computational complexity is $\bigO(n \log n)$ from
the FFT/IFFT rather than $\bigO(n^2)$ from the attention matrix.

\subsection{Why Hardware?}

While spectral mixing can be implemented in software (PyTorch), deploying it as a
custom ASIC provides:
\begin{itemize}[noitemsep]
  \item \textbf{Throughput:} A pipelined FFT engine processes 1 sample/clock
    continuously---no matrix allocation overhead.
  \item \textbf{Power:} Fixed-point arithmetic (Q8.8) at 1.8\,V consumes
    $\sim$4\,mW vs $\sim$30\,W for a GPU running the same computation.
  \item \textbf{Cost:} The chip can be fabricated for \$50--\$500 through Tiny
    Tapeout, vs thousands of dollars for a GPU.
  \item \textbf{Resolution invariance:} Because weights are mode-indexed, the
    same chip processes $n=128$ and $n=512$ sequences without reconfiguration.
\end{itemize}

%═══════════════════════════════════════════════════════════════════════════
\section{Mathematical Foundations}
%═══════════════════════════════════════════════════════════════════════════

\subsection{Self-Attention}

The standard scaled dot-product attention [7] is:
\begin{equation}
  \text{Attn}(Q,K,V) = \text{softmax}\!\left(\frac{QK^\top}{\sqrt{d_k}}\right) V
\end{equation}
The attention matrix $A = QK^\top \in \R^{n \times n}$ requires:
\begin{itemize}[noitemsep]
  \item \textbf{Time:} $\bigO(n^2 d)$ for the matrix multiply + $\bigO(n^2)$ for
    softmax + $\bigO(n^2 d)$ for $AV$.
  \item \textbf{Memory:} $\bigO(n^2)$ for the attention matrix.
  \item \textbf{Parameters:} $\bigO(n^2)$ if position-dependent; $\bigO(d^2)$
    for the $Q,K,V$ projections.
\end{itemize}

\subsection{Fourier Neural Operator (FNO)}

The FNO [1] parameterizes the kernel integral in Fourier space. For
an input $v \in \R^{n \times d}$ (batch of $n$ tokens, $d$ channels):
\begin{equation}
  v_{t+1} = \sigma\!\left( W v_t + \mathcal{K}(v_t) \right)
  \label{eq:fno}
\end{equation}
where $W$ is a linear transform and $\mathcal{K}$ is the spectral kernel:
\begin{align}
  \hat{v} &= \FFT(v) \in \C^{n \times d} \\
  \hat{v}_m' &= R_m \hat{v}_m, \quad m = 0, 1, \ldots, k-1 \\
  \hat{v}_m' &= 0, \quad m = k, \ldots, n-1 \quad \text{(truncation)} \\
  \mathcal{K}(v) &= \IFFT(\hat{v}') \in \R^{n \times d}
\end{align}
Here $R_m \in \C^{d \times d}$ is a learnable complex weight matrix for mode $m$.
In practice, $R$ is truncated to the first $k$ modes (low-frequency components),
which captures the dominant signal structure.

\textbf{Complexity:} The FFT is $\bigO(n \log n)$, the mode-wise multiply is
$\bigO(k d^2)$ (or $\bigO(k d)$ for diagonal $R$), and the IFFT is $\bigO(n \log n)$.
Total: $\bigO(n \log n + k d^2)$.

\subsection{Adaptive Fourier Neural Operator (AFNO)}

The AFNO [2] introduces three modifications to the FNO:

\subsubsection{Block-Diagonal Weight Structure}
Instead of a full $d \times d$ matrix $R_m$ per mode, the channel dimension is split
into $n_b = d / b$ blocks of size $b \times b$:
\begin{equation}
  R_m = \text{diag}(R_m^{(1)}, R_m^{(2)}, \ldots, R_m^{(n_b)})
  \label{eq:blockdiag}
\end{equation}
where each $R_m^{(j)} \in \C^{b \times b}$. This reduces parameters from
$k d^2$ to $k d b$ (a factor of $d/b$ reduction). With $b=1$, the weight is fully
diagonal (per-channel); with $b=d$, it is a single full block (standard FNO).

\subsubsection{Adaptive Soft-Thresholding}
Each spectral coefficient is soft-thresholded to induce sparsity:
\begin{equation}
  \softthr(z, \tau) = \begin{cases}
    0 & \text{if } \abs{z} \leq \tau \\
    z \cdot \frac{\abs{z} - \tau}{\abs{z}} & \text{if } \abs{z} > \tau
  \end{cases}
  \label{eq:soft_threshold}
\end{equation}
where $\tau \geq 0$ is a learnable (or fixed) threshold. This zeros small
coefficients, sparsifying the frequency response. With $\tau = 0$, the AFNO reduces
to a plain FNO.

\subsubsection{Resolution Invariance}
The weight tensor $R$ is indexed by \emph{mode} $m$, not by \emph{position} $i$.
Since the FFT of any length $n \geq k$ produces the same first $k$ modes, the same
weights process sequences of different lengths. Formally:
\begin{proposition}[Resolution Invariance]
\label{prop:resinv}
Let $f_\theta : \R^{n_1 \times d} \to \R^{n_1 \times d}$ be an AFNO layer with
parameters $\theta = \{R_m\}_{m=0}^{k-1}$. Then $f_\theta$ is well-defined for any
input $x \in \R^{n_2 \times d}$ with $n_2 \geq k$, and the parameter count
$\abs{\theta} = k \cdot d \cdot b$ is independent of $n_1$ and $n_2$.
\end{proposition}

\subsection{FFTNet --- Adaptive Spectral Filtering}

FFTNet [3] makes the spectral filter \emph{input-dependent}.
Instead of a static weight $R_m$, the filter is predicted from a global context:
\begin{align}
  g &= \frac{1}{n} \sum_{i=1}^{n} x_i \in \R^d \quad \text{(global context)} \\
  w &= \text{MLP}(g) \in \C^k \quad \text{(predicted filter)} \\
  \hat{v}_m' &= w_m \cdot \hat{v}_m, \quad m = 0, \ldots, k-1
\end{align}
The MLP maps the $d$-dimensional context to $2k$ real values, reshaped into $k$
complex filter coefficients. This adds adaptivity at the cost of the MLP parameters.

\subsection{modReLU Activation}

The modReLU [8] is a magnitude-gated activation for complex
values:
\begin{equation}
  \modReLU(z, b) = z \cdot \sign(\abs{z} + b)
  \label{eq:modrelu}
\end{equation}
where $z \in \C$, $b \in \R$ is a learnable bias, and $\sign(0) = 0$. The bias
shifts the magnitude threshold: positive $b$ makes the activation sparser (more
outputs zeroed); negative $b$ keeps small magnitudes alive.

\subsection{Complexity Comparison}

\begin{table}[h]
\centering
\caption{Complexity comparison: self-attention vs.\ spectral mixing.}
\label{tab:complexity}
\begin{tabular}{lccc}
\toprule
\textbf{Method} & \textbf{Time} & \textbf{Memory} & \textbf{Parameters} \\
\midrule
Self-attention & $\bigO(n^2 d)$ & $\bigO(n^2)$ & $\bigO(n^2)$ \\
FNO (full) & $\bigO(n \log n + k d^2)$ & $\bigO(n d)$ & $\bigO(k d^2)$ \\
AFNO (block-diag) & $\bigO(n \log n + k d b)$ & $\bigO(n d)$ & $\bigO(k d b)$ \\
FFTNet (adaptive) & $\bigO(n \log n + k d)$ & $\bigO(n d)$ & $\bigO(k d + d h)$ \\
\bottomrule
\end{tabular}
\end{table}

For our chip with $n=256$, $d=64$, $k=32$, $b=8$:
\begin{itemize}[noitemsep]
  \item Self-attention: $256^2 \times 64 = 4{,}194{,}304$ MACs
  \item AFNO: $256 \log_2 256 \times 2 + 32 \times 64 \times 8 = 4{,}096 + 16{,}384 = 20{,}480$ ops
  \item Speedup: $\sim$200$\times$
\end{itemize}

%═══════════════════════════════════════════════════════════════════════════
\section{Fixed-Point Arithmetic}
%═══════════════════════════════════════════════════════════════════════════

The chip uses fixed-point arithmetic to avoid the area and power overhead of
floating-point hardware. The primary format is \textbf{Q8.8}: 16-bit total with
8 integer bits (including sign) and 8 fractional bits.

\subsection{Q8.8 Format}

A Q8.8 number represents values in the range $[-128, 127.996]$ with resolution
$2^{-8} = 0.00390625$. The raw integer $r$ maps to the real value:
\begin{equation}
  x = \frac{r}{2^{\text{frac}}}, \quad r \in [-2^{W-1}, 2^{W-1}-1]
  \label{eq:q88}
\end{equation}
where $W = 16$ (total bits), $\text{frac} = 8$. The scale factor is $s = 2^8 = 256$.

\subsection{Saturation Arithmetic}

All arithmetic operations saturate to the representable range:
\begin{equation}
  \text{clamp}(r) = \begin{cases}
    r_{\max} = 2^{W-1} - 1 = 32767 & \text{if } r > r_{\max} \\
    r_{\min} = -2^{W-1} = -32768 & \text{if } r < r_{\min} \\
    r & \text{otherwise}
  \end{cases}
\end{equation}
Overflow and underflow flags are sticky---once set, they remain set until
explicitly cleared. This allows the host to detect precision loss after a
full inference pass.

\subsection{Round-to-Nearest-Even (Banker's Rounding)}

Quantization uses round-half-to-even (banker's rounding) to eliminate
statistical bias:
\begin{equation}
  \text{RNE}(x) = \begin{cases}
    \lfloor x \rfloor & \text{if } x - \lfloor x \rfloor < 0.5 \\
    \lceil x \rceil & \text{if } x - \lfloor x \rfloor > 0.5 \\
    \lfloor x \rfloor & \text{if } x - \lfloor x \rfloor = 0.5 \text{ and } \lfloor x \rfloor \text{ is even} \\
    \lceil x \rceil & \text{if } x - \lfloor x \rfloor = 0.5 \text{ and } \lfloor x \rfloor \text{ is odd}
  \end{cases}
\end{equation}
For integer division by $2^k$, this is implemented as:
\begin{equation}
  \text{rns\_div\_pow2}(v, k) = \sign(v) \cdot \left( \lfloor \frac{\abs{v}}{2^k} \rfloor + \begin{cases} 1 & \text{if remainder} > 2^{k-1} \\ \text{tie-break} & \text{if remainder} = 2^{k-1} \\ 0 & \text{otherwise} \end{cases} \right)
\end{equation}

\subsection{Complex Arithmetic}

\subsubsection{Complex Multiply}
For $z_1 = a + jb$ and $z_2 = c + jd$:
\begin{equation}
  z_1 \cdot z_2 = (ac - bd) + j(ad + bc)
  \label{eq:complex_mul}
\end{equation}
This requires 4 real multiplications and 2 additions. The full-precision product
is $2W$ bits wide; it is rescaled by right-shifting $\text{frac}$ bits (with RNE):
\begin{equation}
  \text{result} = \text{rns\_div\_pow2}(\text{product}, \text{FRAC})
\end{equation}

\subsubsection{Magnitude Approximation}
The exact magnitude $\abs{z} = \sqrt{a^2 + b^2}$ requires a square root, which is
expensive in hardware. The chip uses a shift-and-add approximation:
\begin{equation}
  \abs{z} \approx \max(\abs{a}, \abs{b}) + \frac{1}{2} \min(\abs{a}, \abs{b})
  \label{eq:mag_approx}
\end{equation}
This is accurate to within $\sim$8\% worst case (maximum error at $a = b$, where
the approximation gives $1.5a$ vs.\ the true $\sqrt{2}a \approx 1.414a$, an
error of $\sim$6\%). The $\frac{1}{2}\min$ term is computed by a right shift
(1 bit), requiring no multiplier.

\subsubsection{CORDIC-Lite Magnitude}
An alternative CORDIC-lite algorithm iteratively rotates $(a, b)$ toward the
real axis:
\begin{align}
  a_{i+1} &= a_i - \sigma_i \cdot b_i \cdot 2^{-i} \\
  b_{i+1} &= b_i + \sigma_i \cdot a_i \cdot 2^{-i}
\end{align}
where $\sigma_i = -\sign(b_i)$. After $n$ iterations, $\abs{z} \approx a_n$.
With 8 iterations, accuracy is $\sim$12 bits---sufficient for Q8.8.

%═══════════════════════════════════════════════════════════════════════════
\section{FFT Architecture}
%═══════════════════════════════════════════════════════════════════════════

\subsection{The Discrete Fourier Transform}

The $N$-point DFT of a sequence $x \in \C^N$ is:
\begin{equation}
  X_k = \sum_{n=0}^{N-1} x_n \, e^{-j 2\pi k n / N}, \quad k = 0, \ldots, N-1
  \label{eq:dft}
\end{equation}
Direct computation requires $\bigO(N^2)$ complex multiplications. The FFT reduces
this to $\bigO(N \log N)$ by exploiting the symmetry of the twiddle factors
$W_N^k = e^{-j 2\pi k / N}$.

\subsection{Radix-4 Butterfly}

The chip uses a radix-4 decimation-in-time (DIT) FFT. The radix-4 butterfly takes
4 complex inputs $x_0, x_1, x_2, x_3$ and produces 4 outputs:
\begin{align}
  X_0 &= x_0 + x_1 + x_2 + x_3 \\
  X_1 &= (x_0 - j x_1 - x_2 + j x_3) \cdot W_N^{k_1} \\
  X_2 &= (x_0 - x_1 + x_2 - x_3) \cdot W_N^{k_2} \\
  X_3 &= (x_0 + j x_1 - x_2 - j x_3) \cdot W_N^{k_3}
  \label{eq:radix4}
\end{align}
where multiplication by $j$ swaps real/imaginary parts and negates:
$j(a + jb) = -b + ja$. Each butterfly requires 3 complex twiddle multiplications
(the $X_0$ output has no twiddle).

\subsection{Twiddle Factors}

The twiddle factors are $W_N^k = \cos(2\pi k / N) - j \sin(2\pi k / N)$. For $N=256$,
there are 64 unique twiddle values (stored in a ROM). The symmetry property
$W_N^{k+N/4} = -j \cdot W_N^k$ generates 4 twiddles from 1 stored value, reducing
ROM from 64 to 16 entries (4$\times$ compression).

\subsection{FFT Pipeline for $N=256$}

The 256-point radix-4 FFT has $\log_4 256 = 4$ stages, each with 64 butterflies:
\begin{equation}
  \text{Total butterflies} = \frac{N}{4} \log_4 N = 64 \times 4 = 256
\end{equation}
Total complex multiplies: $3 \times 256 = 768$ (vs.\ $256^2 = 65{,}536$ for direct
DFT). The pipeline processes 1 butterfly per clock cycle, giving a latency of
256 cycles and a throughput of 1 sample/clock.

\subsection{IFFT via Conjugate Method}

The IFFT is computed using the same FFT hardware via the conjugate method:
\begin{equation}
  \IFFT(x) = \frac{1}{N} \conj{\FFT(\conj{x})}
  \label{eq:ifft_conj}
\end{equation}
This requires only input/output conjugation (sign flip of imaginary parts) and
a division by $N$ (right shift by $\log_2 N$ in fixed-point). The V2 architecture
shares a single \texttt{fft\_ifft\_256} instance for both transforms, saving
$\sim$15K gates.

\subsection{Bit-Reversal Permutation}

The DIT FFT produces output in bit-reversed order. For $N=256$, the bit-reversal
permutes 8-bit indices: bit $k \to$ bit $7-k$. This is done in hardware by a
bit-reversal router (a simple crossbar that permutes address bits), saving
$\sim$256 host CPU cycles vs.\ software bit-reversal.

%═══════════════════════════════════════════════════════════════════════════
\section{Chip Architecture}
%═══════════════════════════════════════════════════════════════════════════

\subsection{Pipeline Overview}

The spectral mixer pipeline is:
\begin{equation}
  \text{Input Buffer} \xrightarrow{\text{256 complex}} \text{FFT Engine}
  \xrightarrow{\text{256 modes}} \text{Spectral Multiply}
  \xrightarrow{k=32 \text{ modes}} \text{IFFT Engine}
  \xrightarrow{\text{256 complex}} \text{modReLU}
  \xrightarrow{\text{256 complex}} \text{Output Buffer}
\end{equation}

The chip processes $D=64$ channels sequentially through this pipeline. The
control state machine has states:
\texttt{IDLE} $\to$ \texttt{LOAD\_FFT} $\to$ \texttt{FFT\_RUN} $\to$
\texttt{SM\_IFFT} $\to$ \texttt{MR\_STORE} $\to$ \texttt{DONE}.

\subsection{Parameters}

\begin{table}[h]
\centering
\caption{Chip parameters.}
\label{tab:params}
\begin{tabular}{lc l}
\toprule
\textbf{Parameter} & \textbf{Value} & \textbf{Description} \\
\midrule
$N$ & 256 & FFT size / sequence length \\
$D$ & 64 & Number of channels (serialized) \\
$K$ & 32 & Retained spectral modes \\
$b$ & 8 & Block-diagonal block size \\
$W$ & 16 & Datapath width (Q8.8) \\
$\text{FRAC}$ & 8 & Fractional bits \\
\bottomrule
\end{tabular}
\end{table}

\subsection{Register Map}

The chip exposes 16 32-bit registers via the Wishbone bus:

\begin{table}[h]
\centering
\caption{Register map.}
\label{tab:registers}
\begin{tabular}{lll}
\toprule
\textbf{Address} & \textbf{Name} & \textbf{Description} \\
\midrule
0x00 & CTRL & [0]=start, [1]=done \\
0x04 & STATUS & [0]=busy, [1]=done, [2]=error \\
0x08 & N\_MODES & Number of spectral modes (default 32) \\
0x0C & BLOCK\_SIZE & Block-diagonal block size (default 8) \\
0x10 & THRESHOLD & Soft-thresholding value (Q8.8) \\
0x14 & WEIGHT\_BASE & Weight memory base address \\
0x18 & DATA\_BASE & Data buffer base address \\
0x1C & MODRELU\_BIAS & modReLU bias (Q8.8) \\
0x20 & WEIGHT\_WR\_DATA & Write: re[15:0], im[31:16] \\
0x24 & WEIGHT\_RD\_DATA & Read: re[15:0], im[31:16] \\
0x28 & DATA\_WR\_DATA & Write: re[15:0], im[31:16] \\
0x2C & DATA\_RD\_DATA & Read: re[15:0], im[31:16] \\
0x30 & CONFIG & [7:0]=WIDTH, [15:8]=FRAC, [31:16]=N \\
0x3C & VERSION & Hardware version ID (read-only) \\
\bottomrule
\end{tabular}
\end{table}

\subsection{Wishbone B4 Pipelined Bus}

The V2+ chip uses a Wishbone B4 Pipelined bus interface for 120\,MHz operation:

\begin{table}[h]
\centering
\caption{Wishbone B4 signals.}
\label{tab:wishbone}
\begin{tabular}{ll}
\toprule
\textbf{Signal} & \textbf{Description} \\
\midrule
\texttt{wb\_cyc\_i} & Cycle (bus master requesting) \\
\texttt{wb\_stb\_i} & Strobe (valid request) \\
\texttt{wb\_we\_i} & Write enable \\
\texttt{wb\_adr\_i} & Address (6-bit, word-aligned) \\
\texttt{wb\_dat\_i} & Write data (32-bit) \\
\texttt{wb\_dat\_o} & Read data (32-bit) \\
\texttt{wb\_ack\_o} & Acknowledge (1-cycle delayed in pipelined mode) \\
\texttt{wb\_stall\_o} & Stall (chip not ready for next request) \\
\texttt{wb\_bte\_i} & Burst type extension (00=single, 01=4-word, etc.) \\
\texttt{wb\_cti\_i} & Cycle type (000=classic, 010=incr burst, 111=end) \\
\bottomrule
\end{tabular}
\end{table}

The pipelined ack breaks the combinational path at 120\,MHz (8.3\,ns cycle):
the request is latched on cycle $N$, and the acknowledge is emitted on cycle
$N+1$, giving a full clock period for address decode and register access.

\subsubsection{Burst Mode}
Burst transactions auto-increment the address for consecutive weight or data
writes. With \texttt{wb\_cti\_i} = \texttt{010} (incrementing burst), the
internal \texttt{burst\_addr} register auto-increments, and the host does not need
to drive a new address each cycle. When \texttt{wb\_cti\_i} = \texttt{111} (end of
burst), the burst terminates after the current word. This reduces bus overhead
from 1 cycle/word to $\sim$0.25 cycles/word for 4-word bursts.

\subsection{Gate Count Breakdown}

\begin{table}[h]
\centering
\caption{V1 gate count estimate.}
\label{tab:gates_v1}
\begin{tabular}{lr}
\toprule
\textbf{Module} & \textbf{Gates} \\
\midrule
FFT engine (radix-4, 256-point) & $\sim$15{,}000 \\
Spectral weight multiply (32 modes) & $\sim$8{,}000 \\
IFFT engine (shared with FFT) & $\sim$15{,}000 (shared) \\
modReLU + soft-threshold & $\sim$2{,}000 \\
Wishbone interface + control & $\sim$500 \\
\midrule
\textbf{Total} & $\sim$40{,}000 \\
\bottomrule
\end{tabular}
\end{table}

%═══════════════════════════════════════════════════════════════════════════
\section{Quantization}
%═══════════════════════════════════════════════════════════════════════════

\subsection{Symmetric int8 Quantization}

Complex spectral weights are quantized from float32 to int8 using symmetric
quantization (zero-point = 0):
\begin{equation}
  q = \text{clip}\!\left( \text{round}\!\left( \frac{x}{s} \right), -128, 127 \right), \quad
  s = \frac{\max(\abs{x})}{127}
  \label{eq:quant}
\end{equation}
Dequantization is: $x \approx q \cdot s$. The scale $s$ is stored per-layer
(separate scales for real and imaginary parts).

\subsection{Quantization Error}

The quantization error for a single value is bounded by:
\begin{equation}
  \abs{x - q \cdot s} \leq \frac{s}{2} = \frac{\max(\abs{x})}{254}
\end{equation}
For typical spectral weights with $\max(\abs{x}) \sim 1.0$, the error is
$\sim 0.004$ (0.4\%), well within the $\sim$2\% threshold measured in practice.

\subsection{Int8 FFT}

The chip includes an int8 FFT implementation using lookup tables for twiddle
factors. The twiddle factors are pre-quantized to int8 and stored in a ROM.
Products are computed in int32 (to prevent overflow) and rescaled by dividing
by $q_{\max} = 127$:
\begin{equation}
  \text{prod} = (w_r \cdot b_r - w_i \cdot b_i) \, // \, q_{\max}
\end{equation}
Per-stage scaling (right-shift by 1 every other stage) prevents overflow,
giving a total scaling of $\sim 1/\sqrt{N}$.

%═══════════════════════════════════════════════════════════════════════════
\section{Resolution Invariance}
%═══════════════════════════════════════════════════════════════════════════

\begin{theorem}[Resolution Invariance]
\label{thm:resinv}
Let $f_\theta$ be a spectral mixing layer (FNO/AFNO) with parameters
$\theta = \{R_m\}_{m=0}^{k-1}$. Then:
\begin{enumerate}[noitemsep]
  \item $f_\theta$ accepts inputs of any sequence length $n \geq k$.
  \item The number of parameters $\abs{\theta} = k \cdot d \cdot b$ is
    independent of $n$.
  \item The same trained parameters work at different $n$ without
    retraining.
\end{enumerate}
\end{theorem}

\begin{proof}
The forward pass is:
\begin{equation}
  f_\theta(x) = x + \IFFT_n\!\left( R \cdot \FFT_n(x) \right)
\end{equation}
where $\FFT_n$ and $\IFFT_n$ are $n$-point transforms. The weight $R$ operates
on the first $k$ modes of $\FFT_n(x)$. Since the first $k$ Fourier modes of an
$n$-point FFT correspond to the same frequencies regardless of $n$ (they are
$e^{-j2\pi m / N}$ for $m = 0, \ldots, k-1$), the weight $R_m$ applies to the same
frequency component at any $n \geq k$. Thus:
\begin{itemize}[noitemsep]
  \item The weight tensor has shape $(k, d, b)$, which does not depend on $n$.
  \item The FFT/IFFT are non-parametric operations whose size is determined at
    runtime by $x.\text{shape}[1]$.
  \item No retraining is needed because the weights operate on frequency modes,
    not spatial positions.
\end{itemize}
\end{proof}

\textbf{Practical implication:} A model trained at $n=64$ can be evaluated at
$n=256$ or $n=512$ with zero parameter changes---only the FFT size changes.
This is impossible with standard attention, which requires $n \times n$ weight
matrices that must be retrained for each $n$.

%═══════════════════════════════════════════════════════════════════════════
\section{V2 Architectural Improvements}
%═══════════════════════════════════════════════════════════════════════════

The V2 architecture introduces three key improvements over V1:

\subsection{Shared FFT/IFFT Engine (Improvement 1)}

V1 uses separate \texttt{fft\_256} and \texttt{ifft\_256} modules ($\sim$15K gates
each). V2 replaces them with a single \texttt{fft\_ifft\_256} module that uses
the conjugate method (Eq.~\ref{eq:ifft_conj}) with a mode bit selecting FFT vs.
IFFT. Since the FFT and IFFT run sequentially (FFT finishes before IFFT starts),
time-multiplexing is possible. This saves $\sim$15K gates.

\subsection{Merged Soft-Threshold + modReLU (Improvement 8)}

V1 applies soft-thresholding (in spectral multiply) and modReLU (separate module)
in two stages with two separate magnitude computations. V2 fuses them:
\begin{equation}
  \text{merged}(z) = \begin{cases}
    0 & \text{if } \abs{z} < \tau \;\;(\text{soft-threshold}) \\
    0 & \text{if } \abs{z} + b \leq 0 \;\;(\text{modReLU}) \\
    z & \text{otherwise}
  \end{cases}
\end{equation}
Both checks use the same $\abs{z}$ (computed once), saving one magnitude unit
and one pipeline stage ($\sim$1K gates).

\subsection{Serialized Channel Processing (Improvement 9)}

V1 processes $D=64$ channels in parallel (64 spectral multiply units). V2
serializes: one channel at a time through a single MAC unit, with a channel
counter. This trades $8\times$ latency for $8\times$ smaller spectral multiply
area. For LLM inference (batch=1), this is an acceptable trade-off.

\begin{table}[h]
\centering
\caption{V1 vs.\ V2 comparison.}
\label{tab:v1v2}
\begin{tabular}{lrrr}
\toprule
\textbf{Metric} & \textbf{V1} & \textbf{V2} & \textbf{Change} \\
\midrule
Total area (GE) & 100{,}300 & 9{,}102 & $-90.9\%$ \\
Total power (mW) & 32.50 & 2.21 & $-93.2\%$ \\
Throughput (tokens/s) & 189{,}394 & 2{,}959 & $0.016\times$ \\
Latency (cycles, $k=32$) & 544 & 34{,}816 & $64\times$ \\
\bottomrule
\end{tabular}
\end{table}

The V2 area reduction comes from: shared FFT engine ($-15$K), serialized
channels ($-76$K), merged activation ($-1$K), and removed IFFT module ($-15$K).

%═══════════════════════════════════════════════════════════════════════════
\section{V3 Performance Improvements}
%═══════════════════════════════════════════════════════════════════════════

V3 adds 20 improvements organized into four groups:

\subsection{Datapath Arithmetic (1--5)}

\subsubsection{Booth-Encoded Radix-4 Multiplier}
Booth encoding [9] halves the number of partial products
in the complex multiplier. For a radix-4 Booth encoding:
\begin{equation}
  \text{partial products} = \frac{W}{2} \quad \text{(vs.\ $W$ for standard)}
\end{equation}
Carry-save (Wallace tree) compression eliminates the carry-propagate adder
from the critical path, reducing it by $\sim$30\% and enabling 50$\to$65\,MHz.

\subsubsection{Block Floating-Point (BFP)}
Each block of 64 FFT samples shares a 4-bit exponent with 12-bit mantissas:
\begin{equation}
  x = m \cdot 2^e, \quad e = \lfloor \log_2 \max(\abs{x_i}) \rfloor
\end{equation}
This gives $\sim$12 extra bits of dynamic range (162.5\,dB vs.\ 90.3\,dB for
Q8.8), reducing FFT rounding error by $\sim$20$\times$.

\subsubsection{Carry-Save Accumulator}
Partial products are kept in carry-save format throughout $k$-mode
accumulation, with a single carry-propagate at the end:
\begin{equation}
  \text{latency} = 1 \;\text{cycle} \quad \text{(vs.\ $k$ cycles for carry-propagate each step)}
\end{equation}
83\% latency reduction in the MAC accumulator.

\subsubsection{FMA Butterfly}
The twiddle multiply and butterfly add/sub are fused:
\begin{equation}
  \text{FMA}(a, W, b) = a + W \cdot b \quad \text{(single operation, no intermediate rounding)}
\end{equation}
Saves 1 pipeline stage and $\sim$500 gates per butterfly, with improved numerical
accuracy from eliminating intermediate rounding.

\subsubsection{Truncated Booth Multiplier}
For twiddle factors in $[-1, 1]$, the integer part of the product is bounded.
Computing only the lower 16 bits of a $16 \times 16$ product saves $\sim$30\%
of the multiplier area with negligible error ($\sim$0.004$).

\subsection{Memory and Data Movement (6--10)}

\subsubsection{Ping-Pong Dual Buffers}
Two memory banks alternate between FFT stages (A$\to$B, B$\to$A), eliminating
pipeline bubbles for 6.8$\times$ throughput improvement.

\subsubsection{Shadow Weight Prefetch}
A shadow register file prefetches the next weight block while the current
block is processed. Latency reduction: 48.4\%, with zero stalls when load
time $\leq$ compute time.

\subsubsection{Zero-Skip MAC with Dummy Cycles}
Zeroed modes (from soft-thresholding) use dummy cycles with LFSR decoy data:
\begin{equation}
  \text{if } \abs{z_m} < \tau: \;\text{compute decoy}(LFSR_m) \text{ instead of } z_m \cdot w_m
\end{equation}
Constant timing preserved, 45\% power reduction at 50\% sparsity. An attacker
cannot distinguish real from dummy multiplies via timing or power.

\subsubsection{Conflict-Free Addressing}
Modulo-4 bank mapping: $\text{bank} = \text{addr}[1:0]$, $\text{row} = \text{addr} \gg 2$.
Guarantees butterfly data in different banks $\to$ 0 stalls, single-cycle execution.

\subsubsection{Bit-Reversal Router}
Hardware crossbar that permutes address bits ($\text{bit}[k] \to \text{bit}[\log_2 N - 1 - k]$),
saving 255 CPU cycles vs.\ software bit-reversal.

\subsection{Algorithmic Optimizations (11--15)}

\subsubsection{Real-Input FFT (RFFT)}
For real-valued inputs, Hermitian symmetry gives $N/2+1$ modes instead of $N$:
\begin{equation}
  X_{N-k} = \conj{X_k} \quad \Rightarrow \quad \text{only } X_0, \ldots, X_{N/2} \text{ are unique}
\end{equation}
49.6\% mode reduction (129 vs.\ 256 modes for $N=256$).

\subsubsection{Twiddle Factor Symmetry}
$W_N^{k+N/4} = -j \cdot W_N^k$ generates 4 twiddles from 1 stored value:
$W^k$, $-jW^k$, $-W^k$, $jW^k$. 4$\times$ ROM compression (64$\to$16 entries).

\subsubsection{Mode Interleaving}
Even/odd modes processed simultaneously in 2 pipeline stages:
$2\times$ spectral multiply throughput without $2\times$ area.

\subsubsection{Adaptive Mode Count}
Configurable $k$ (8--32) via Wishbone register. Fewer modes = faster inference.
Timing varies with $k$ (public config), not with data content.

\subsubsection{Early IFFT Start}
IFFT begins after $k$ modes are ready (not all $N$), overlapping with remaining
FFT output. 5.9\% latency reduction (512$\to$480 cycles for $k=32$).

\subsection{Pipeline and Throughput (16--20)}

\subsubsection{Deep 8-Stage Pipeline}
2 sub-stages per radix-4 stage (twiddle multiply + butterfly add). Shorter
critical path enables 65$\to$80\,MHz (60\% frequency boost).

\subsubsection{Dual-Channel Parallel}
Two independent spectral channels share weight storage: $2\times$ throughput
for batch-2 inference at $2\times$ FFT area.

\subsubsection{Configurable FFT Size}
128/256/512-point FFT via Wishbone register. $2\times$ speedup for short
sequences; $2\times$ longer context for long ones. Same weights work at any size
(resolution invariance).

\subsubsection{DVFS with Secure Tracking}
Dynamic voltage-frequency scaling (1.8V$\to$1.2V) with secure voltage tracker.
60\% idle power reduction, 78\% active power reduction at low $k$. Transitions
only between batches (constant intra-batch timing).

\subsubsection{DMA Burst Controller}
4-word burst transfers via pipelined Wishbone B3: 75\% bus overhead reduction
(1$\to$0.25 cycles/word). Weight load: 128$\to$32 cycles.

%═══════════════════════════════════════════════════════════════════════════
\section{Security Architecture}
%═══════════════════════════════════════════════════════════════════════════

The chip implements 11 security measures, all preserved across V1--V3:

\begin{table}[h]
\centering
\caption{Security measures.}
\label{tab:security}
\begin{tabular}{ll}
\toprule
\textbf{Measure} & \textbf{Mechanism} \\
\midrule
Bitstream encryption & LFSR-based stream cipher on weight bitstream \\
Constant-time MAC & All $k$ modes processed in fixed cycles regardless of data \\
Power flattening & Decoy MAC with LFSR data masks real multiply activity \\
Integrity hash & 32-bit rolling hash of weight memory \\
EM shielding & Top metal layer EM shield over compute core \\
Reproducible build & Manifest hash verification of all source files \\
DVFS secure tracking & Voltage transition verified before frequency change \\
\bottomrule
\end{tabular}
\end{table}

\subsection{Constant-Time Spectral MAC}

The constant-time guarantee ensures that the number of clock cycles for spectral
multiplication depends only on $k$ (the configured mode count), not on the input
data:
\begin{equation}
  T_{\text{MAC}} = k \cdot T_{\text{mode}} \quad \text{(independent of input values)}
\end{equation}
Zeroed modes (from soft-thresholding) are processed through dummy cycles with
LFSR-generated decoy data, so the cycle count and power profile are identical
whether a mode is zeroed or not.

\subsection{Power Flattening}

The decoy MAC computes meaningless multiplies on LFSR-generated data with the
same statistical properties as real data. The power consumption of a real
multiply and a dummy multiply are indistinguishable:
\begin{equation}
  P_{\text{real}}(w, x) \approx P_{\text{decoy}}(\text{LFSR}, \text{LFSR}) \quad \forall w, x
\end{equation}
This prevents power-analysis side-channel attacks from recovering spectral
weight values.

\subsection{Integrity Hash}

A 32-bit rolling hash covers the weight memory:
\begin{equation}
  H = \sum_{i=0}^{N-1} w_i \cdot 2^{i \bmod 32} \pmod{2^{32}}
\end{equation}
The hash is computed on-chip at initialization and compared against the expected
value loaded by the host. Any tampering with weights (bit-flip, substitution)
is detected with probability $1 - 2^{-32}$.

%═══════════════════════════════════════════════════════════════════════════
\section{Benchmark Results}
%═══════════════════════════════════════════════════════════════════════════

\subsection{Area Comparison}

\begin{table}[h]
\centering
\caption{Area comparison across versions (gate equivalents).}
\label{tab:area}
\begin{tabular}{lrrr}
\toprule
\textbf{Module} & \textbf{V1} & \textbf{V2} & \textbf{V3} \\
\midrule
FFT engine & 5{,}100 & 5{,}100 & 4{,}870 \\
IFFT engine & 5{,}100 & 0 (shared) & 0 (shared) \\
Spectral multiply & 76{,}800 & 1{,}200 & 800 \\
modReLU & 12{,}800 & 200 & 200 \\
Wishbone/control & 500 & 500 & 800 \\
Security modules & 0 & 2{,}000 & 2{,}100 \\
Ping-pong buffers & 0 & 0 & 8{,}000 \\
Clock gating & 0 & 102 & 121 \\
DVFS controller & 0 & 0 & 300 \\
Deep pipeline regs & 0 & 0 & 400 \\
Shadow prefetch & 0 & 0 & 500 \\
Bit-rev router & 0 & 0 & 300 \\
Conflict-free addr & 0 & 0 & 200 \\
\midrule
\textbf{Total} & \textbf{100{,}300} & \textbf{9{,}102} & \textbf{18{,}591} \\
Reduction vs.\ V1 & --- & 90.9\% & 81.5\% \\
\bottomrule
\end{tabular}
\end{table}

\subsection{Power Comparison}

\begin{table}[h]
\centering
\caption{Power comparison across versions (mW).}
\label{tab:power}
\begin{tabular}{lrrr}
\toprule
\textbf{Module} & \textbf{V1} & \textbf{V2} & \textbf{V3} \\
\midrule
FFT engine & 1.800 & 1.800 & 4.280 \\
IFFT engine & 1.800 & 0 & 0 \\
Spectral multiply & 25.600 & 0.400 & 0.140 \\
modReLU & 3.200 & 0.050 & 0.050 \\
Wishbone/control & 0.100 & 0.100 & 0.100 \\
Security (decoy MAC) & 0 & 0.400 & 0.400 \\
Clock gating savings & 0 & $-0.540$ & $-0.642$ \\
DVFS savings & 0 & 0 & $-0.428$ \\
Deep pipeline regs & 0 & 0 & 0.150 \\
\midrule
\textbf{Total} & \textbf{32.50} & \textbf{2.21} & \textbf{4.05} \\
Reduction vs.\ V1 & --- & 93.2\% & 87.5\% \\
\bottomrule
\end{tabular}
\end{table}

\subsection{Throughput Comparison}

\begin{table}[h]
\centering
\caption{Throughput comparison (tokens/s) at various $k$ and $N$.}
\label{tab:throughput}
\begin{tabular}{lrrr}
\toprule
\textbf{Config} & \textbf{V1} & \textbf{V2} & \textbf{V3} \\
\midrule
$k=8$, $N=128$ & 189{,}394 & 2{,}959 & 19{,}841 \\
$k=8$, $N=256$ & 96{,}154 & 1{,}502 & 9{,}843 \\
$k=8$, $N=512$ & 48{,}450 & 757 & 4{,}902 \\
$k=16$, $N=128$ & 183{,}824 & 2{,}872 & 20{,}492 \\
$k=16$, $N=256$ & 94{,}697 & 1{,}480 & 10{,}000 \\
$k=24$, $N=128$ & 178{,}571 & 2{,}790 & 21{,}186 \\
$k=32$, $N=128$ & 173{,}611 & 2{,}713 & 21{,}930 \\
$k=32$, $N=256$ & 91{,}912 & 1{,}436 & 10{,}331 \\
\midrule
\textbf{Max} & \textbf{189{,}394} & \textbf{2{,}959} & \textbf{21{,}930} \\
V3 improvement vs.\ V2 & --- & --- & $7.41\times$ \\
\bottomrule
\end{tabular}
\end{table}

\subsection{Latency Comparison}

\begin{table}[h]
\centering
\caption{Latency comparison ($k=32$, $N=256$).}
\label{tab:latency}
\begin{tabular}{lrrr}
\toprule
\textbf{Stage} & \textbf{V1} & \textbf{V2} & \textbf{V3} \\
\midrule
FFT cycles & 256 & 256 & 129 \\
Spectral multiply cycles & 32 & 32 & 16 \\
IFFT cycles & 256 & 256 & 256 \\
\midrule
\textbf{Total} & \textbf{544} & \textbf{34{,}816} & \textbf{7{,}744} \\
Change vs.\ V1 & --- & $+6300\%$ & $+1324\%$ \\
\bottomrule
\end{tabular}
\end{table}

V1's low latency is due to 64$\times$ parallel channel processing (expensive in
area). V2 serializes channels (64$\times$ latency increase). V3 recovers via
dual-channel, deep pipeline, RFFT, and early IFFT.

%═══════════════════════════════════════════════════════════════════════════
\section{Silicon Target}
%═══════════════════════════════════════════════════════════════════════════

\subsection{SkyWater SKY130 Process}

\begin{table}[h]
\centering
\caption{SKY130 process parameters.}
\label{tab:sky130}
\begin{tabular}{ll}
\toprule
\textbf{Parameter} & \textbf{Value} \\
\midrule
Node & 130\,nm \\
Metal layers & 6 (1P6M) \\
Core voltage & 1.8\,V \\
I/O voltage & 3.3\,V \\
Standard cell density & $\sim$100K routed gates/mm$^2$ \\
SRAM available & 1KB--64KB macros \\
\bottomrule
\end{tabular}
\end{table}

\subsection{Tiny Tapeout}

Tiny Tapeout provides the cheapest path to silicon:
\begin{itemize}[noitemsep]
  \item Cost: \$50--\$500 per tile
  \item Area: $\sim$1\,mm$^2$ ($\sim$100K gates)
  \item I/O: Limited (shared pins, SPI-like protocol)
  \item Timeline: $\sim$3 months from submission to packaged chip
\end{itemize}

The spectral mixer core ($\sim$40K gates V1, $\sim$18.6K gates V3) fits within a
single Tiny Tapeout tile.

\subsection{ChipIgnite Scaling}

ChipIgnite ($\sim$\$14{,}950) provides $\sim$10\,mm$^2$ for a full transformer:
\begin{table}[h]
\centering
\caption{ChipIgnite area budget.}
\label{tab:chipignite}
\begin{tabular}{lr}
\toprule
\textbf{Component} & \textbf{Area (mm$^2$)} \\
\midrule
Spectral mixer core (V3) & 0.19 \\
4 transformer layers (spectral + FFN) & 3.5 \\
Token embeddings (16KB SRAM) & 0.10 \\
Unembedding (16KB SRAM) & 0.10 \\
FFN weights (32KB SRAM) & 0.20 \\
\midrule
\textbf{Total} & $\sim$4.1 \\
Available & 10.0 \\
\bottomrule
\end{tabular}
\end{table}

\subsection{Design Flow}

\begin{equation}
  \text{Verilog RTL} \xrightarrow{\text{Yosys}} \text{Gate netlist}
  \xrightarrow{\text{OpenROAD}} \text{GDSII}
  \xrightarrow{\text{DRC/LVS}} \text{Tapeout}
\end{equation}

The open-source EDA flow: Yosys (synthesis) $\to$ OpenROAD (floorplan, placement,
CTS, routing) $\to$ OpenSTA (timing) $\to$ Magic + Netgen (DRC/LVS) $\to$ GDSII.

%═══════════════════════════════════════════════════════════════════════════
\section{Compiler and Software Stack}
%═══════════════════════════════════════════════════════════════════════════

\subsection{SpectralChipCompiler}

The compiler converts a trained PyTorch spectral transformer into a binary blob:
\begin{enumerate}[noitemsep]
  \item Walk the model and discover all spectral layers (AFNO/FFTNet).
  \item Extract complex weight tensors $R_m \in \C^{k \times d \times b}$.
  \item Quantize real and imaginary parts separately to int8 (symmetric, zero-point=0).
  \item Pack into a binary blob with header (magic \texttt{SPLR}, version, metadata).
\end{enumerate}

\subsubsection{Binary Format}
\begin{verbatim}
Header (64 bytes):
  magic        : 4s  "SPLR"
  version      : I   1
  n_layers     : I
  d_model      : I
  seq_len      : I
  n_modes      : I
  block_size   : I

Per-layer block:
  layer_type   : I   0=AFNO, 1=FFTNet
  weight_len   : I
  scale_real   : f   float32
  scale_imag   : f   float32
  threshold    : f   float32
  weights_real : [int8] * weight_len
  weights_imag : [int8] * weight_len
\end{verbatim}

\subsection{Python Simulation Stack}

\begin{table}[h]
\centering
\caption{Python modules.}
\label{tab:python}
\begin{tabular}{ll}
\toprule
\textbf{Module} & \textbf{Purpose} \\
\midrule
\texttt{fno.py} & Fourier Neural Operator layer \\
\texttt{afno.py} & Adaptive FNO with soft-thresholding \\
\texttt{fftnet.py} & FFTNet adaptive spectral filter \\
\texttt{transformer.py} & Spectral transformer block \\
\texttt{model.py} & Tiny spectral LM ($\sim$100K params) \\
\texttt{fixedpoint.py} & Q8.8 arithmetic simulator \\
\texttt{quantize.py} & int8 quantization + int8 FFT \\
\texttt{compiler.py} & PyTorch $\to$ chip blob compiler \\
\texttt{perf\_sim.py} & V3 performance simulators (20 modules) \\
\texttt{security.py} & Integrity hash + LFSR cipher \\
\texttt{constant\_time.py} & Constant-time spectral MAC \\
\bottomrule
\end{tabular}
\end{table}

\subsection{Test Suite}

The project has 343 passing tests (11 slow tests deselected):
\begin{itemize}[noitemsep]
  \item \texttt{test\_afno.py}: AFNO reduction to FNO, gradients, thresholds, block sizes
  \item \texttt{test\_fno.py}: Shape, zero modes, gradients, resolution invariance
  \item \texttt{test\_fftnet.py}: modReLU, forward pass, sequence lengths
  \item \texttt{test\_fixedpoint.py}: Q-format arithmetic, overflow, complex multiply
  \item \texttt{test\_quantize.py}: Quantization, dequantization, int8 FFT
  \item \texttt{test\_compiler.py}: Compile/load round-trip, verification
  \item \texttt{test\_complexity.py}: Complexity benchmark
  \item \texttt{test\_transformer.py}: Forward pass, gradients, resolution invariance
  \item \texttt{test\_perf\_sim.py}: All 20 V3 improvement simulators
  \item \texttt{test\_security.py}: Integrity hash, bitstream encryption, tamper detection
  \item \texttt{test\_constant\_time.py}: Constant-timing verification
  \item \texttt{test\_resolution.py}: Resolution invariance tests
\end{itemize}

%═══════════════════════════════════════════════════════════════════════════
\section{V4--V7 Additional Improvements}
%═══════════════════════════════════════════════════════════════════════════

Beyond V3's 20 improvements, the project adds 33 more across V4--V7:

\subsection{V4 Speed Improvements (21--25)}
\begin{itemize}[noitemsep]
  \item \textbf{Triple Twiddle ROM}: 3 parallel read ports, 3$\times$ faster twiddle fetch
  \item \textbf{Streaming IFFT Loader}: Overlapped IFFT input with spectral multiply output
  \item \textbf{Mode-Skip Multiply}: Skip multiplier for truncated modes ($\geq K$), 87.5\% power savings
  \item \textbf{Pipelined Butterfly}: 2-stage multiply/add split, 40\% critical path reduction, 65$\to$90\,MHz
  \item \textbf{Batch Channel Controller}: Auto-sequencing for all $D=64$ channels, 6{,}300 cycles saved/layer
\end{itemize}

\subsection{V6 Clock Speed Improvements (26--32)}
\begin{itemize}[noitemsep]
  \item \textbf{Retimed Pipeline Registers}: Balanced path delays, 100$\to$120\,MHz
  \item \textbf{Clock Tree Optimizer}: H-tree with 8 buffer stages, $<50$\,ps skew
  \item \textbf{Multi-Stage Spectral Multiply}: 3-stage pipeline, 90$\to$110\,MHz
  \item \textbf{Weight Register Banking}: 4-bank parallel read, 0 stall cycles
  \item \textbf{Operand Isolation}: Clock gating inactive stages, 25\% power reduction
  \item \textbf{Wider Datapath (Q12.4)}: 16$\times$ more dynamic range, prevents multi-layer overflow
  \item \textbf{STA Fixup Buffers}: Buffer chains on long paths, 90$\to$120+\,MHz
\end{itemize}

\subsection{V6 Architecture (33--40)}
\begin{itemize}[noitemsep]
  \item \textbf{Speculative IFFT}: Starts at 80\% modes ready, 4\% latency reduction
  \item \textbf{Multi-Port Weight SRAM}: Dual-port read, halves spectral multiply phase
  \item \textbf{Fused FFT+IFFT}: Shared entire datapath, saves $\sim$8K gates
  \item \textbf{Channel Interleaving}: 2$\times$ autoregressive throughput
  \item \textbf{Configurable Pipeline Depth}: 2/4/8 stages, optimal latency-frequency tradeoff
  \item \textbf{Wishbone Burst Write}: 4$\times$ faster input data loading
  \item \textbf{Output Streaming FIFO}: Backpressure, 87\% output wait reduction
  \item \textbf{Layer Scheduler}: Auto weight swapping, 400 cycles saved for 4-layer model
\end{itemize}

\subsection{V6 Reliability (41--45)}
\begin{itemize}[noitemsep]
  \item \textbf{Parity Error Detection}: 1-bit parity per butterfly, single-bit error detection
  \item \textbf{Guard Bands}: Saturate at 95\% of max, prevents cascading overflow
  \item \textbf{Sticky Overflow Counter}: 16-bit total overflow count per pass
  \item \textbf{Redundant Checksum}: 16-bit XOR hash of spectral multiply outputs
  \item \textbf{Thermal Throttle}: Ring-oscillator sensor, graceful $k$ reduction at 75$^\circ$C
\end{itemize}

\subsection{V7 Advanced Butterfly Cores (46--53)}
\begin{itemize}[noitemsep]
  \item \textbf{Split-Radix}: 35\% fewer multiplies (498 vs.\ 768 for $N=256$)
  \item \textbf{Radix-8}: 33\% fewer multiplies, fewer stages (3 vs.\ 4)
  \item \textbf{4$\times$ Parallel Array}: 4 simultaneous butterflies, 4$\times$ throughput
  \item \textbf{Constant-Geometry FFT}: No bit-reversal needed, saves 256 cycles
  \item \textbf{Stockham FFT}: Natural-order output, eliminates bit-reversal hardware
  \item \textbf{Fused Butterfly+Multiply}: FMA, 30\% critical path reduction
  \item \textbf{CORDIC Twiddle}: On-the-fly generation, saves $\sim$1{,}200 gates (ROM-free)
  \item \textbf{Radix-8 Twiddle Bank}: 7 simultaneous twiddle factors, 7$\times$ faster access
\end{itemize}

Total across all versions: \textbf{53 improvements}.

%═══════════════════════════════════════════════════════════════════════════
\section{Conclusion and Future Work}
%═══════════════════════════════════════════════════════════════════════════

\subsection{Summary of Findings}

Spectral Silicon demonstrates that the $\bigO(n^2)$ self-attention bottleneck in
LLM inference can be replaced by an $\bigO(n \log n)$ spectral mixing pipeline
implemented directly in silicon on a 130\,nm process. Key findings:

\begin{enumerate}
  \item \textbf{Area efficiency}: V3 achieves 18{,}591 GE---an 81.5\% reduction
    from V1's 100{,}300 GE---while maintaining all 11 security measures.
  \item \textbf{Power efficiency}: V3 consumes 4.05\,mW (87.5\% reduction from
    V1's 32.5\,mW), making it feasible for battery-powered edge deployment.
  \item \textbf{Throughput}: V3 achieves 21{,}930 tokens/s (7.41$\times$ over V2),
    enabled by dual-channel processing, mode interleaving, and an 80\,MHz clock.
  \item \textbf{Resolution invariance}: The same trained weights process
    sequences of different lengths ($n=128, 256, 512$) without retraining---a
    property unique to spectral mixing and impossible with self-attention.
  \item \textbf{Quantization tolerance}: int8 symmetric quantization introduces
    $<2\%$ error, compatible with the Q8.8 fixed-point datapath.
  \item \textbf{Security}: Constant-time MAC with dummy cycles and power
    flattening via decoy multiplies are \emph{combined} with power reduction
    (zero-skip MAC), demonstrating that security and efficiency can coexist.
\end{enumerate}

\subsection{Future Work}

\begin{itemize}
  \item \textbf{Tapeout}: Submit V3 or V6 to a Tiny Tapeout shuttle for
    fabrication and silicon validation.
  \item \textbf{Analog extension}: Explore charge-domain FFT (switched-capacitor)
    on SKY130 analog tiles for ultra-low-power spectral mixing.
  \item \textbf{Multi-layer SoC}: Use ChipIgnite's 10\,mm$^2$ to integrate
    4 complete spectral transformer layers with embeddings, FFN, and sampling.
  \item \textbf{Larger FFT}: Support $N=512$ or $N=1024$ for longer context
    windows, leveraging the resolution-invariance property.
  \item \textbf{Training study}: Compare AFNO vs.\ FFTNet vs.\ standard attention
    on downstream NLP tasks (perplexity, BLEU) at matched parameter counts.
  \item \textbf{Feed-forward network chip}: A companion ASIC for the FFN stage
    of the transformer, completing the on-chip inference pipeline.
\end{itemize}

%═══════════════════════════════════════════════════════════════════════════
\section*{References}
%═══════════════════════════════════════════════════════════════════════════

\addcontentsline{toc}{section}{References}

\begin{enumerate}[label={[\arabic*]}]
  \item Li, Z.\ et al. ``Fourier Neural Operator for Parametric PDEs.''
    \emph{ICLR} 2021. arXiv:2010.08895.
  \item Guibas, J.\ et al. ``Adaptive Fourier Neural Operators: Efficient
    Token Mixers for Transformers.'' 2021. arXiv:2111.13587.
  \item Fein-Ashley, J.\ ``The FFT Strikes Back: An Efficient Alternative to
    Self-Attention.'' 2025. arXiv:2502.18394.
  \item Lee-Thorp, J.\ et al. ``FNet: Mixing Tokens with Fourier Transforms.''
    \emph{NeurIPS} 2021.
  \item Lu, L.\ et al. ``Learning Nonlinear Operators via DeepONet.''
    \emph{Nature Machine Intelligence}, 2021.
  \item Kovachki, N.\ et al. ``Neural Operator: Learning Maps Between
    Complex Function Spaces.'' 2021. arXiv:2108.08481.
  \item Vaswani, A.\ et al. ``Attention Is All You Need.''
    \emph{NeurIPS} 2017.
  \item Arjovsky, M.\ et al. ``Unitary Evolution Recurrent Neural Networks.''
    \emph{ICML} 2016. (modReLU definition)
  \item Booth, A.D.\ ``A Signed Binary Multiplication Technique.''
    \emph{Quarterly Journal of Mechanics and Applied Mathematics}, 1951.
  \item SkyWater SKY130 PDK: \url{https://github.com/google/skywater-pdk}
  \item Tiny Tapeout: \url{https://tinytapeout.com/}
  \item Efabless ChipIgnite: \url{https://chipfoundry.io/}
  \item OpenROAD Project: \url{https://theopenroadproject.org/}
  \item OpenLane: \url{https://github.com/The-OpenROAD-Project/OpenLane}
\end{enumerate}

\end{document}
